\uuid{ZL9o}
\titre{Base orthonormée}
\auteur{bodin ; hanani ; mzari}
\organisation{exo7}
\licence{CC BY-NC-SA 4.0}
\datecreate{2020-01-20}
\niveau{L1}
\module{Géométrie}
\chapitre{Géométrie du plan}
\sousChapitre{Distance, norme, produit scalaire}
\difficulte{3}
\theme{base orthonormée}
\exo7qcmid{lille-1.geometrie-plane.021}

\contenu{
\begin{question}
\qclassification{140.01,140.02}
\qtags{motcle=base orthonormée}
Dans le plan muni d'un repère orthonormé direct $(O,\vec{i},\vec{j})$, on considère les vecteurs $\displaystyle \vec{u}=\left(\frac{1}{2},a\right)$ et $\displaystyle \vec{v}=\left(-\frac{\sqrt{3}}{2},b\right)$. Comment choisir les réels $a$ et $b$ pour que $(\vec{u},\vec{v})$ soit une base orthonormée directe ?
\begin{answers}
    \good{$\displaystyle a=\frac{\sqrt{3}}{2}$ et $\displaystyle b=\frac{1}{2}$}
    \bad{$\displaystyle a=\frac{\sqrt{3}}{2}$ et $\displaystyle b=-\frac{1}{2}$}
    \bad{$\displaystyle a=-\frac{\sqrt{3}}{2}$ et $\displaystyle b=\frac{1}{2}$}
    \bad{$\displaystyle a=-\frac{\sqrt{3}}{2}$ et $\displaystyle b=-\frac{1}{2}$}
\end{answers}
\begin{explanations}
D'abord, $\displaystyle \|\vec{u}\|=1\Leftrightarrow a=\frac{\pm \sqrt{3}}{2}$, $\displaystyle \|\vec{v}\|=1\Leftrightarrow b=\frac{\pm 1}{2}$ et $\vec{u}\cdot\vec{v}=0$ si, et seulement si, $a$ et $b$ sont de même signe. Enfin pour que $(\vec{u},\vec{v})$ soit directe, il faut que $\mbox{det}(\vec{u},\vec{v})$ soit positif.
\end{explanations}
\end{question}
}
